% GENERATED FILE. Edit canonical lessons, then run npm run generate:books.
% canonical-source-hash: fnv1a64:6c9e894804e447c2
% canonical-lessons: BN-C115-bash, BN-C115-amgachh, BN-C115-narkel, BN-C115-anaros, BN-C115-pepe

\chapter{Bamboo, Mango tree, Coconut, Pineapple, Papaya}
\label{ch:bn-bamboo-mango-tree-coconut-pineapple-papaya}
\hlchaptermodality{115}

\begin{chapteropening}
\textbf{By the end of this chapter:} \emph{I can say bamboo, a mango tree, a coconut, a pineapple and a papaya.}

Complete the last lesson of chapter 115: I can say bamboo, a mango tree, a coconut, a pineapple and a papaya.
\end{chapteropening}

\section[\texorpdfstring{\bn{বাঁশ} (bā̃sh)}{বাঁশ (bā̃sh)}]{\bn{বাঁশ} (bā̃sh) — bamboo}

\label{lesson:BN-C115-bash}



\begin{quote}
Before the new one: say the Bengali for a seed, then the Bengali for a thorn.
\end{quote}

\subsection*{You'll want to know: \bn{বাঁশ}}
\textbf{\bn{বাঁশ}} — \emph{bā̃sh} — \textquotedblleft{}bamboo\textquotedblright{}.

\begin{cousinweb}
\textbf{\bn{বাঁশ}} (\emph{bāṁsh}) is a tadbhava, a word inherited from Sanskrit and reshaped by sound change. It comes from \emph{vaṁśa}, \textquotedblleft{}bamboo\textquotedblright{} (CDIAL 11175).
\end{cousinweb}

\begin{grammarlens}[title={What you've built}]
One more word for this chapter.
\end{grammarlens}

\subsection*{Your turn}
\begin{itemize}
\raggedright
  \item \emph{Say it:} \emph{bā̃sh}
  \item \emph{Say it:} \emph{bā̃sh}, once more
  \item \emph{Say it:} say \emph{bā̃sh}
  \item \emph{Recall:} say the Bengali for a seed, then the Bengali for a thorn, then say \emph{bā̃sh} again
\end{itemize}

\begin{tcolorbox}[breakable,colback=teal!4,colframe=teal!35!black,title={Before you move on}]
What does \bn{বাঁশ} mean? (\textquotedblleft{}Bamboo\textquotedblright{}.) Say it once more.
\end{tcolorbox}
\section[\texorpdfstring{\bn{আমগাছ} (āmgāchh)}{আমগাছ (āmgāchh)}]{\bn{আমগাছ} (āmgāchh) — a mango tree}

\label{lesson:BN-C115-amgachh}



\begin{quote}
Before the new one: say the Bengali for a thorn, then the Bengali for bamboo.
\end{quote}

\subsection*{You'll want to know: \bn{আমগাছ}}
\textbf{\bn{আমগাছ}} — \emph{āmgāchh} — \textquotedblleft{}a mango tree\textquotedblright{}.

\begin{cousinweb}
\textbf{\bn{আমগাছ}} (\emph{āmgāchh}) is \emph{ām}, \textquotedblleft{}mango\textquotedblright{}, from Sanskrit \emph{āmra} (CDIAL 1268), + \emph{gāchh}, \textquotedblleft{}tree\textquotedblright{}, from \emph{*gakṣa}, a form CDIAL reconstructs (CDIAL 3949).
\end{cousinweb}

\begin{grammarlens}[title={What you've built}]
One more word for this chapter.
\end{grammarlens}

\subsection*{Your turn}
\begin{itemize}
\raggedright
  \item \emph{Say it:} \emph{āmgāchh}
  \item \emph{Say it:} \emph{āmgāchh}, once more
  \item \emph{Say it:} say \emph{āmgāchh}
  \item \emph{Recall:} say the Bengali for a thorn, then the Bengali for bamboo, then say \emph{āmgāchh} again
\end{itemize}

\begin{tcolorbox}[breakable,colback=teal!4,colframe=teal!35!black,title={Before you move on}]
What does \bn{আমগাছ} mean? (\textquotedblleft{}A mango tree\textquotedblright{}.) Say it once more.
\end{tcolorbox}
\section[\texorpdfstring{\bn{নারকেল} (nārkel)}{নারকেল (nārkel)}]{\bn{নারকেল} (nārkel) — a coconut}

\label{lesson:BN-C115-narkel}



\begin{quote}
Before the new one: say the Bengali for bamboo, then the Bengali for a mango tree.
\end{quote}

\subsection*{You'll want to know: \bn{নারকেল}}
\textbf{\bn{নারকেল}} — \emph{nārkel} — \textquotedblleft{}a coconut\textquotedblright{}.

\begin{cousinweb}
\textbf{\bn{নারকেল}} (\emph{nārkel}) goes back to Sanskrit \emph{nārikela}, \textquotedblleft{}coconut palm and fruit\textquotedblright{}. CDIAL gives no Bengali form under that entry, but it lists Oriya \emph{naṛhera} and \emph{naṛiā}, \textquotedblleft{}coconut\textquotedblright{}, the forms nearest to Bengali (CDIAL 7075). CDIAL reports Bloch's view that \emph{nārikela} is a Dravidian compound.
\end{cousinweb}

\begin{grammarlens}[title={What you've built}]
One more word for this chapter.
\end{grammarlens}

\subsection*{Your turn}
\begin{itemize}
\raggedright
  \item \emph{Say it:} \emph{nārkel}
  \item \emph{Say it:} \emph{nārkel}, once more
  \item \emph{Say it:} say \emph{nārkel}
  \item \emph{Recall:} say the Bengali for bamboo, then the Bengali for a mango tree, then say \emph{nārkel} again
\end{itemize}

\begin{tcolorbox}[breakable,colback=teal!4,colframe=teal!35!black,title={Before you move on}]
What does \bn{নারকেল} mean? (\textquotedblleft{}A coconut\textquotedblright{}.) Say it once more.
\end{tcolorbox}
\section[\texorpdfstring{\bn{আনারস} (ānārôs)}{আনারস (ānārôs)}]{\bn{আনারস} (ānārôs) — a pineapple}

\label{lesson:BN-C115-anaros}



\begin{quote}
Before the new one: say the Bengali for a mango tree, then the Bengali for a coconut.
\end{quote}

\subsection*{You'll want to know: \bn{আনারস}}
\textbf{\bn{আনারস}} — \emph{ānārôs} — \textquotedblleft{}a pineapple\textquotedblright{}.

\begin{cousinweb}
\textbf{\bn{আনারস}} (\emph{ānārôs}) is a loan from Portuguese \emph{ananás}, \textquotedblleft{}pineapple\textquotedblright{}.
\end{cousinweb}

\begin{grammarlens}[title={What you've built}]
One more word for this chapter.
\end{grammarlens}

\subsection*{Your turn}
\begin{itemize}
\raggedright
  \item \emph{Say it:} \emph{ānārôs}
  \item \emph{Say it:} \emph{ānārôs}, once more
  \item \emph{Say it:} say \emph{ānārôs}
  \item \emph{Recall:} say the Bengali for a mango tree, then the Bengali for a coconut, then say \emph{ānārôs} again
\end{itemize}

\begin{tcolorbox}[breakable,colback=teal!4,colframe=teal!35!black,title={Before you move on}]
What does \bn{আনারস} mean? (\textquotedblleft{}A pineapple\textquotedblright{}.) Say it once more.
\end{tcolorbox}
\section[\texorpdfstring{\bn{পেঁপে} (pepe)}{পেঁপে (pepe)}]{\bn{পেঁপে} (pepe) — a papaya}

\label{lesson:BN-C115-pepe}



\begin{quote}
Before the new one: say the Bengali for a coconut, then the Bengali for a pineapple.
\end{quote}

\subsection*{You'll want to know: \bn{পেঁপে}}
\textbf{\bn{পেঁপে}} — \emph{pepe} — \textquotedblleft{}a papaya\textquotedblright{}.

\begin{cousinweb}
\textbf{\bn{পেঁপে}} (\emph{pepe}) is a loan from Portuguese \emph{papaia}, \textquotedblleft{}papaya\textquotedblright{}. Portuguese took the word from a Caribbean language.
\end{cousinweb}

\begin{grammarlens}[title={What you've built}]
One more word for this chapter.
\end{grammarlens}

\subsection*{Your turn}
\begin{itemize}
\raggedright
  \item \emph{Say it:} \emph{pepe}
  \item \emph{Say it:} \emph{pepe}, once more
  \item \emph{Say it:} say \emph{pepe}
  \item \emph{Recall:} say the Bengali for a coconut, then the Bengali for a pineapple, then say \emph{pepe} again
\end{itemize}

\begin{tcolorbox}[breakable,colback=teal!4,colframe=teal!35!black,title={Before you move on}]
What does \bn{পেঁপে} mean? (\textquotedblleft{}A papaya\textquotedblright{}.) Say it once more.
\end{tcolorbox}
